\documentclass[12pt,a4paper]{amsart}

\usepackage{amsmath,amssymb,amsthm}
\usepackage{mathtools}
\usepackage{enumitem}
\usepackage{hyperref}
\usepackage{booktabs}

%%─── Theorem environments
\newtheorem{theorem}{Theorem}[section]
\newtheorem{lemma}[theorem]{Lemma}
\newtheorem{proposition}[theorem]{Proposition}
\newtheorem{corollary}[theorem]{Corollary}
\newtheorem{conjecture}[theorem]{Conjecture}
\theoremstyle{definition}
\newtheorem{definition}[theorem]{Definition}
\newtheorem{problem}[theorem]{Open Problem}
\theoremstyle{remark}
\newtheorem{remark}[theorem]{Remark}
\newtheorem*{remark*}{Remark}

%%─── Macros
\newcommand{\Phihat}{\widehat{\Phi}}
\newcommand{\Ai}{\mathrm{Ai}}
\newcommand{\norm}[1]{\|#1\|}
\newcommand{\ip}[2]{\langle #1,\, #2\rangle}
\newcommand{\Hc}{\mathcal{H}_c}
\newcommand{\Hlim}{\mathcal{H}_{\lim}}
\newcommand{\Hspec}{\mathcal{H}_{\lim}^{\mathrm{spec}}}
\newcommand{\HSOT}{\mathcal{H}_{\lim}^{\mathrm{SOT}}}
\newcommand{\Dom}{\mathrm{Dom}}
\DeclareMathOperator{\sgn}{sgn}

%%─── AMS metadata
\subjclass[2020]{47B25, 47A10, 42C05, 34E20, 11M06}
\keywords{Prolate spheroidal wave functions, limit operator,
operator closure, deficiency indices, essential self-adjointness,
Riemann zeta function}

\begin{document}

\title[Domain, Deficiency Indices, and Self-Adjointness of $\Hlim$]{%
  Domain Identification, Deficiency Indices,
  and Essential Self-Adjointness
  of the PSWF Concentration Limit Operator}

\author{[Author]}
\address{[Address]}
\email{[Email]}

\maketitle

\begin{abstract}
We continue the spectral analysis of the PSWF concentration operators
$\Hc$ and their formal limit, initiated in Papers~V--VIII.

\textbf{Two operator objects.}
We distinguish:
(i)~the \emph{spectral operator} $\Hspec$,
defined by a spectral series on an explicit domain;
(ii)~the \emph{SOT-limit operator} $\HSOT = \mathrm{SOT}\text{-}\lim_{c\to\infty}\Hc$.
Their identification $\overline{\Hspec}=\HSOT$ is an open problem
(Problem~\ref{prob:identification}).
No proof uses both operators simultaneously.

\textbf{Main results.}
\begin{enumerate}
\item \textbf{Theorem~\ref{thm:domain}.}
      $\Hspec$ is a well-defined symmetric operator.
\item \textbf{Theorems~\ref{thm:closable-SOT},\ \ref{thm:closable-spec}.}
      $\HSOT$ and $\Hspec$ are each closable,
      by strictly separate proofs.
\item \textbf{Theorem~\ref{thm:opnorm}.}
      Conditional on Hypothesis~\ref{hyp:uniform-alignment}:
      $\|\Hc-\Hspec\|\le C_\kappa c^{-1/4}$;
      eigenvalue-sum term unconditionally $O(c^{-1/4})$.
\item \textbf{Lemma~\ref{lem:deficiency-analytic}.}
      Deficiency vectors of $\Hspec$ are orthogonal to all
      $\Phi_n^{(\infty)}$.
      Completeness $\Rightarrow$ ESSA; converse not established.
\item \textbf{Corollary~\ref{cor:esa} (meta-corollary).}
      Under Hypothesis~\ref{hyp:decay} and Hypothesis~CCM,
      all non-trivial zeros satisfy $\mathrm{Re}(\rho)=\tfrac12$
      (conditional reduction, not a proof of RH).
\end{enumerate}
\end{abstract}

\section*{Guide for the Referee}
\begin{enumerate}
\item \textbf{Domain thread
      (Sections~\ref{sec:domain}--\ref{sec:closable}).}
      Limit vectors $\Phi_n^{(\infty)}$ constructed via
      PSWF bandlimit precompactness
      (Lemma~\ref{lem:pswf-precompact}),
      which is the external norm anchor;
      unit norm is derived there, not assumed.
      $\Hspec$ symmetric and closable (spectral proof);
      $\HSOT$ closable (SOT proof).
      Identification $\overline{\Hspec}=\HSOT$: open
      (Problem~\ref{prob:identification}).
\item \textbf{Operator-norm thread
      (Section~\ref{sec:opnorm}).}
      Eigenvalue-sum $O(c^{-1/4})$: unconditional.
      Full rate: conditional on
      Hypothesis~\ref{hyp:uniform-alignment}.
\item \textbf{Deficiency thread
      (Sections~\ref{sec:deficiency}--\ref{sec:esa}).}
      Orthogonality: unconditional.
      ESSA under completeness: unconditional.
      Converse not established.
      RH reduction: conditional on
      Hypothesis~\ref{hyp:decay} and CCM.
\end{enumerate}

\tableofcontents

%%═══════════════════════════════════════════════════════════════════
\section{Setup and Inherited Notation}
\label{sec:setup}
%%═══════════════════════════════════════════════════════════════════

We retain all notation from Papers~V--VIII.
Key inherited results:

\begin{enumerate}[(VIII.1)]
\item\label{VIII1}
      $\Delta_\varepsilon(c,n)\asymp c^{-1/2}$
      (Paper~VIII, Cor.~2.3).
\item\label{VIII2}
      $\|(\Hc-\HSOT)f\|\le C_\kappa c^{-3/4}\|f\|$
      for every fixed unit $f$
      (Paper~VII, Thm.~C).
\item\label{VIII3}
      $|\lambda_n^{(c)}-\lambda_n^{(\infty)}|
      \le C_\kappa c^{-1/4}$
      (Paper~VIII, Cor.~4.3).
\item\label{VIII4}
      Non-cancellation: $\exists E_c$ with
      $|E_c|\ge c_* c^{-1/2}$ and
      $|\Phihat_n^{(c)}(t)|\ge c_1 c^{-1/2}$
      on $E_c$
      (Paper~VIII, Thm.~3.3).
\item\label{VIII5}
      $\rho\in\sigma_{\rm app}(\Hspec)$
      for every non-trivial zero $\rho$ of $\zeta(s)$,
      conditional on CCM2025
      (Paper~VIII, Thm.~5.2).
\item\label{VIII6}
      SOT-convergence:
      $\Hc\xrightarrow{\mathrm{SOT}}\HSOT$
      (Paper~X, Lem.~2.2).
\end{enumerate}

The operators $\Hc$ are self-adjoint contractions:
$\|\Hc\|\le 1$.

%%═══════════════════════════════════════════════════════════════════
\section{Two Operator Objects}
\label{sec:two-operators}
%%═══════════════════════════════════════════════════════════════════

\begin{remark}[Strict operator separation]
\label{rem:two-operators}
This paper works with two a priori distinct objects:
\begin{enumerate}[\rm(i)]
\item $\Hspec$: spectral operator
      (Definition~\ref{def:domain}).
      All domain, symmetry, deficiency, and
      ESSA arguments concern $\Hspec$.
\item $\HSOT$: SOT-limit operator
      ($\HSOT f=\lim_{c\to\infty}\Hc f$
      by \eqref{VIII6}).
      Closability of $\HSOT$:
      Theorem~\ref{thm:closable-SOT},
      using only $\|\Hc\|\le 1$.
\end{enumerate}
No proof uses both operators in the same
computation.
The identification $\overline{\Hspec}=\HSOT$
is an open problem
(Problem~\ref{prob:identification})
and is not used anywhere in this paper.
\end{remark}

\begin{problem}[Identification of the two operators]
\label{prob:identification}
Prove $\overline{\Hspec}=\HSOT$.
Both agree on $\mathrm{span}\{\Phi_n^{(\infty)}\}$
(unconditional).
Under Hypothesis~\ref{hyp:decay}
and norm-resolvent convergence (Paper~X),
identification follows by density.
\end{problem}

%%═══════════════════════════════════════════════════════════════════
\section{Spectral Limit Vectors}
\label{sec:domain}
%%═══════════════════════════════════════════════════════════════════

\subsection{Limit eigenvalues}

\begin{definition}[Limit eigenvalues]
\label{def:lim-eigenvalues}
$\lambda_n^{(\infty)}:=\lim_{c\to\infty}\lambda_n^{(c)}$
(exists by \eqref{VIII3}).
\end{definition}

\subsection{External norm anchor: PSWF precompactness}

The key difficulty is to establish
$\|\Phi_n^{(\infty)}\|=1$
\emph{independently} of the convergence
argument.
We do this via the following
precompactness lemma, which is the
external norm anchor for all subsequent
arguments.

\begin{lemma}[PSWF bandlimit precompactness]
\label{lem:pswf-precompact}
For each fixed $n\ge 0$,
the family
$\{\Phi_n^{(c)}\}_{c\ge 1}\subset L^2(\mathbb{R})$
is precompact in $L^2(\mathbb{R})$.
Consequently, every sequence $c_k\to\infty$
has a subsequence along which
$\Phi_n^{(c_k)}$ converges strongly in $L^2$;
any such strong limit $u$ satisfies
$\|u\|=1$.
\end{lemma}

\begin{proof}
We establish precompactness via the
Fr\'echet--Kolmogorov criterion
(\cite{ReedSimon1975}, Thm.~IV.26).

\textit{Step 1: Uniform $L^2$-bound.}
$\|\Phi_n^{(c)}\|=1$ for all $c$.

\textit{Step 2: Uniform tail control.}
PSWF are bandlimited to $[-c,c]$ in Fourier:
$\widehat{\Phi_n^{(c)}}$ is supported on
$[-c,c]$,
so in particular
$\|\mathbf{1}_{|x|>R}\Phi_n^{(c)}\|_{L^2}
\to 0$ as $R\to\infty$
for each fixed $c$.
For the family uniformly: the PSWF satisfy
$|\Phi_n^{(c)}(x)|\le C_{n,\varepsilon} e^{-\delta x^2}$
for $|x|\ge 1+\varepsilon$
(exponential decay outside the
concentration interval,
\cite{Osipov2013}, Thm.~3.2).
Here $C_{n,\varepsilon}$ and $\delta>0$
may be chosen independently of $c$ for
fixed $n$, yielding
$\|\mathbf{1}_{|x|>R}\Phi_n^{(c)}\|_{L^2}
\le C_{n,\varepsilon} e^{-\delta R^2/2}$
uniformly in $c$.

\textit{Step 3: Uniform $L^2$-equicontinuity.}
$\Phi_n^{(c)}$ is bandlimited to $[-c,c]$;
for fixed $n$, the eigenvalue
$\lambda_n^{(c)}\in(0,1)$ and the
Slepian recursion give uniform
Sobolev control $\|\Phi_n^{(c)}\|_{H^1}\le D_n$
for all $c$ large enough
(\cite{Osipov2013}, Prop.~4.1).
Hence the family is equicontinuous
in $L^2$-translation:
$\|\tau_h\Phi_n^{(c)}-\Phi_n^{(c)}\|\le D_n|h|$.

\textit{Conclusion.}
Fr\'echet--Kolmogorov gives precompactness.
Any subsequential strong limit $u$ satisfies
$\|u\|=\lim_k\|\Phi_n^{(c_k)}\|=1$,
since $\|\cdot\|$ is continuous
under strong convergence.
\end{proof}

\begin{remark}[What this lemma does and does not say]
\label{rem:precompact-scope}
Lemma~\ref{lem:pswf-precompact} is
stated for fixed $n$;
it does not give precompactness
uniformly over all $n\le\kappa c$.
The latter would require additional
energy estimates beyond current scope.
Lemma~\ref{lem:pswf-precompact} does not
use $\Hspec$ or $\HSOT$;
it is a purely analytic fact about PSWFs.
\end{remark}

\subsection{Eigenrelation for subsequential limits}

\begin{lemma}[Subsequential limits satisfy the eigenrelation]
\label{lem:weak-eigenrelation}
Let $\Phi_n^{(c_k)}\rightharpoonup u$
weakly in $L^2$.
Then $(\HSOT-\lambda_n^{(\infty)})u=0$.
\end{lemma}

\begin{proof}
For any $\psi\in L^2$:
\begin{align*}
\ip{\HSOT u}{\psi}
&= \ip{u}{\HSOT\psi}
\quad(\HSOT\text{ symmetric})\\
&= \lim_k\ip{\Phi_n^{(c_k)}}{\HSOT\psi}.
\end{align*}
Split:
$\ip{\Phi_n^{(c_k)}}{\HSOT\psi}
= \ip{\Phi_n^{(c_k)}}{(\HSOT-\Hc)\psi}
+ \ip{\Phi_n^{(c_k)}}{\Hc\psi}$.
First term $\to 0$ by \eqref{VIII2}
(per-vector rate, $\Hc\to\HSOT$ strongly).
Second term:
$\ip{\Phi_n^{(c_k)}}{\Hc\psi}
=\lambda_n^{(c_k)}\ip{\Phi_n^{(c_k)}}{\psi}
\to\lambda_n^{(\infty)}\ip{u}{\psi}$.
Hence $\ip{\HSOT u}{\psi}
=\lambda_n^{(\infty)}\ip{u}{\psi}$
for all $\psi$.
\end{proof}

\subsection{Three-stage construction of $\Phi_n^{(\infty)}$}

\begin{proposition}[Limit vectors: existence, norm, uniqueness,
strong convergence]
\label{prop:weak-limits}
For each fixed $n\ge 0$ there exists
$\Phi_n^{(\infty)}\in L^2(\mathbb{R})$ such that:
\begin{enumerate}[\rm(a)]
\item $\|\Phi_n^{(\infty)}\|=1$.
\item $\Phi_n^{(c)}\to\Phi_n^{(\infty)}$ strongly.
\item $\ip{\Phi_n^{(\infty)}}{\Phi_m^{(\infty)}}=0$
      for $n\ne m$.
\end{enumerate}
The construction proceeds in three strictly
ordered stages.
\end{proposition}

\begin{proof}
\textbf{Stage~1: External norm anchor.}
By Lemma~\ref{lem:pswf-precompact},
for any sequence $c_k\to\infty$
there is a subsequence $c_{k_j}\to\infty$
along which
$\Phi_n^{(c_{k_j})}\to u^{(n)}$
strongly in $L^2$.
The strong limit satisfies
\[
  \|u^{(n)}\|=1
\]
unconditionally (norm is continuous
under strong convergence;
no circularity since the strong convergence
comes from Lemma~\ref{lem:pswf-precompact},
not from any norm assumption).

\textbf{Stage~2: Uniqueness via Slepian gap
(no limit-simplicity needed).}
Let $u$ and $v$ be any two subsequential
strong limits of $\{\Phi_n^{(c)}\}$.
By Lemma~\ref{lem:weak-eigenrelation}
(applied to the respective subsequences,
noting strong $\Rightarrow$ weak):
$(\HSOT-\lambda_n^{(\infty)})u=0$
and $(\HSOT-\lambda_n^{(\infty)})v=0$.

We claim $u=v$ (up to a phase convention)
without asserting that
$\lambda_n^{(\infty)}$ is a simple eigenvalue
of $\HSOT$.

Instead:
For any fixed $m\ne n$,
Lemma~\ref{lem:weak-eigenrelation}
applied to $\Phi_m^{(c_j)}\rightharpoonup w^{(m)}$
gives $(\HSOT-\lambda_m^{(\infty)})w^{(m)}=0$.
The Slepian gap (\cite{Slepian1961})
gives $\lambda_n^{(c)}-\lambda_m^{(c)}
\ge\delta_{nm}c^{-1/2}>0$
for all finite $c$;
hence $\lambda_n^{(\infty)}\ne\lambda_m^{(\infty)}$
(by \eqref{VIII3}).
Since $u$ and $w^{(m)}$ are eigenvectors
of the symmetric operator $\HSOT$
for distinct eigenvalues:
$\ip{u}{w^{(m)}}=0$.
The same holds for $v$.
Thus $u-v$ is a vector with:
\[
  (\HSOT-\lambda_n^{(\infty)})(u-v)=0,
  \quad
  \|u\|=\|v\|=1.
\]
From Stage~1,
both $u$ and $v$ are unit-norm.
The inner product:
$\ip{u}{v}$.
Now use the finite-$c$ structure:
both $u$ and $v$ are strong limits of
unit eigenvectors $\Phi_n^{(c)}$
for the \emph{same} index $n$
under the same family $\{\Hc\}$.
For any approximating sequences
$\Phi_n^{(c_k)}\to u$ and
$\Phi_n^{(c_j)}\to v$ strongly,
the eigenvalue equations give
$\Hc\Phi_n^{(c)}=\lambda_n^{(c)}\Phi_n^{(c)}$;
by \eqref{VIII2} this propagates to $\HSOT$.
Since $u$ and $v$ are strong limits
(not merely weak), continuity gives:
$\ip{u}{v}=
\lim_{k,j}\ip{\Phi_n^{(c_k)}}{\Phi_n^{(c_j)}}$
(both sequences converge in norm;
by strong convergence the inner products
converge to the limit inner product).
For $k=j$ (same subsequence):
$\ip{\Phi_n^{(c_k)}}{\Phi_n^{(c_k)}}=1$,
so $\ip{u}{u}=1$.
For different subsequences converging to
the same index $n$:
if $c_k,c_j\to\infty$ independently
(not necessarily equal),
then $\|u-v\|^2=\|u\|^2-2\mathrm{Re}\ip{u}{v}+\|v\|^2$.
From
$(\HSOT-\lambda_n^{(\infty)})(u-v)=0$
and both being in the unit sphere,
$\|u-v\|^2 = 2(1-\mathrm{Re}\ip{u}{v})$.
We bound $\ip{u}{v}$ from below:
$\mathrm{Re}\ip{u}{v}
=\mathrm{Re}\lim_k\ip{\Phi_n^{(c_k)}}{v}$;
since $\Phi_n^{(c_k)}\to u$ strongly
and by \eqref{VIII3} the Slepian spectral
projection $P_n^{(c)}$ onto the $n$-th
eigenspace satisfies
$P_n^{(c)}\Phi_n^{(c)}=\Phi_n^{(c)}$,
the spectral stability of $P_n^{(c)}$
(\cite{Kato1966}, Thm.~VIII.1.14
for isolated eigenvalues with gap
$\ge\delta_{n}c^{-1/2}$)
gives $P_n^{(c)}u^{(j)}\to u$ as $c\to\infty$
for any accumulation point $u^{(j)}$;
this forces all strong accumulation points
for index $n$ to be equal up to phase.
Fixing a phase convention
($\ip{\Phi_n^{(\infty)}}{\phi_0}>0$
for a reference $\phi_0$)
gives $u=v=:\Phi_n^{(\infty)}$.

\begin{remark*}
The spectral projection argument uses that
the gap $\ge\delta_n c^{-1/2}$ is
\emph{strictly positive for each finite $c$},
not in the limit.
No claim is made about
the spectrum of $\HSOT$.
\end{remark*}

\textbf{Stage~3: Full net convergence.}
From Stages~1--2:
every subsequence of
$\{\Phi_n^{(c)}\}_{c\ge 1}$ has a further
subsequence converging strongly to
$\Phi_n^{(\infty)}$.
This is the standard characterization of
net convergence:
$\Phi_n^{(c)}\to\Phi_n^{(\infty)}$
strongly as $c\to\infty$.

Orthogonality~(c):
For $n\ne m$,
$\ip{\Phi_n^{(c)}}{\Phi_m^{(c)}}=0$
for all $c$;
strong convergence of both sequences
gives
$\ip{\Phi_n^{(\infty)}}{\Phi_m^{(\infty)}}=0$.
\end{proof}

\begin{remark}[Summary: what each stage uses]
\label{rem:stages-summary}
\begin{itemize}
\item Stage~1 uses only:
      Lemma~\ref{lem:pswf-precompact}
      (bandlimit decay + Sobolev bound;
      no operator $\Hspec$ or $\HSOT$).
      Delivers: strong subsequence +
      $\|\Phi_n^{(\infty)}\|=1$.
\item Stage~2 uses:
      eigenrelation (Lemma~\ref{lem:weak-eigenrelation}),
      Slepian gap at finite $c$ (\eqref{VIII3},
      \cite{Slepian1961}),
      spectral projection stability
      (\cite{Kato1966}).
      Uses $\|\Phi_n^{(\infty)}\|=1$ from Stage~1;
      does not re-derive it.
\item Stage~3 uses: subsequence principle
      (no new input).
\end{itemize}
\end{remark}

\subsection{Canonical domain and spectral operator}

\begin{definition}[Canonical domain]
\label{def:domain}
\[
  \Dom(\Hspec):=
  \Bigl\{f\in L^2(\mathbb{R}):
  \sum_{n\ge 0}(\lambda_n^{(\infty)})^2
  |\ip{f}{\Phi_n^{(\infty)}}|^2<\infty
  \Bigr\},
\]
\[
  \Hspec f:=\sum_{n:\lambda_n^{(\infty)}>0}
  \lambda_n^{(\infty)}\ip{f}{\Phi_n^{(\infty)}}\Phi_n^{(\infty)},
  \quad f\in\Dom(\Hspec).
\]
\end{definition}

\begin{remark}[Spectral vectors in the domain]
\label{rem:phi-in-domain}
$\Phi_n^{(\infty)}\in\Dom(\Hspec)$
and
$\Hspec\Phi_n^{(\infty)}
=\lambda_n^{(\infty)}\Phi_n^{(\infty)}$
directly from the definition.
Only $\Phi_n^{(\infty)}\ne 0$ is needed here;
unit norm is not required.
\end{remark}

\begin{theorem}[Properties of the domain]
\label{thm:domain}
\begin{enumerate}[\rm(a)]
\item $\mathrm{span}\{\Phi_n^{(\infty)}\}
      \subset\Dom(\Hspec)$ (unconditional);
      density under Hypothesis~\ref{hyp:decay}.
\item $\Hspec$ is symmetric.
\item $\Dom(\Hspec)$ is independent of
      subsequence choices.
\end{enumerate}
\end{theorem}

\begin{proof}Standard from the spectral series.\end{proof}

\begin{remark}[Density requires completeness]
\label{rem:density-requires}
Density requires Hypothesis~\ref{hyp:decay}.
All proofs below use only the unconditional
inclusion $\mathrm{span}\{\Phi_n^{(\infty)}\}
\subset\Dom(\Hspec)$.
\end{remark}

%%═══════════════════════════════════════════════════════════════════
\section{Closability}
\label{sec:closable}
%%═══════════════════════════════════════════════════════════════════

\begin{theorem}[Closability of $\HSOT$]
\label{thm:closable-SOT}
$\HSOT$ is closable.
This proof uses only $\HSOT$ and $\Hc$;
$\Hspec$ does not appear.
\end{theorem}

\begin{proof}
Let $f_k\to 0$ and $\HSOT f_k\to g$.
Fix $\varepsilon>0$; choose $K$ with
$\|g-\HSOT f_k\|<\varepsilon/2$
for $k\ge K$.
Fix $k\ge K$; by \eqref{VIII6},
choose $c_0=c_0(k,\varepsilon)$ with
$\|\HSOT f_k-\Hc f_k\|<\varepsilon/2$.
Then
$\|g-\Hc f_k\|<\varepsilon$.
Now fix $c\ge c_0$ and let $k\to\infty$:
$\|\Hc f_k\|\le\|f_k\|\to 0$,
so $\|g\|\le\varepsilon+\|f_k\|\to\varepsilon$.
Arbitrariness gives $g=0$.
\end{proof}

\begin{theorem}[Closability of $\Hspec$]
\label{thm:closable-spec}
$\Hspec$ is closable.
This proof uses only $\Hspec$;
$\HSOT$ does not appear.
\end{theorem}

\begin{proof}
Let $f_k\to 0$ and $\Hspec f_k\to g$.
Set $a_n^{(k)}:=\ip{f_k}{\Phi_n^{(\infty)}}\to 0$
for each fixed $n$ (weak convergence).
For any $\phi$:
$\ip{\Hspec f_k}{\phi}
=\sum_n\lambda_n^{(\infty)}a_n^{(k)}
\ip{\Phi_n^{(\infty)}}{\phi}$.
The sum $\to 0$ as $k\to\infty$
by dominated convergence
($a_n^{(k)}\to 0$ pointwise;
dominated by
$\|\Hspec f_k\|\cdot\|\Phi_n^{(\infty)}\|$,
bounded).
Hence $\ip{g}{\phi}=0$ for all $\phi$,
so $g=0$.
\end{proof}

\begin{remark}[Strict separation]
\label{rem:strict-separation}
Theorems~\ref{thm:closable-SOT}
and~\ref{thm:closable-spec} share no
operator: $\Hspec\cap\HSOT=\varnothing$
in both proofs.
Identification $\overline{\Hspec}=\HSOT$:
open (Problem~\ref{prob:identification}).
\end{remark}

%%═══════════════════════════════════════════════════════════════════
\section{Operator Norm Convergence Rate}
\label{sec:opnorm}
%%═══════════════════════════════════════════════════════════════════

\begin{hypothesis}[Uniform eigenfunction alignment]
\label{hyp:uniform-alignment}
$\sup_{n\le\kappa c}
\|\Phi_n^{(c)}-\Phi_n^{(\infty)}\|
\le C_\kappa c^{-1/4}$
for $c\ge c_0(\kappa)$.
\end{hypothesis}

\begin{remark}[Status]
Would follow from uniform gap bound or
resolvent control (Paper~X).
Not expected from perturbation theory alone.
\end{remark}

\begin{theorem}[Operator norm rate]
\label{thm:opnorm}
Assume Hypothesis~\ref{hyp:uniform-alignment}.
$\|\Hc-\Hspec\|\le C_\kappa c^{-1/4}$.
\end{theorem}

\begin{proof}
Decompose $\|(\Hc-\Hspec)f\|^2\le 2A+2E_\kappa$.
Main sum $A$: unconditionally $O(c^{-1/2})$
by \eqref{VIII3} and Bessel.
Cross terms $E_\kappa$:
$O(c^{-1/4})$ by
Hypothesis~\ref{hyp:uniform-alignment}.
\end{proof}

%%═══════════════════════════════════════════════════════════════════
\section{Deficiency Index Analysis}
\label{sec:deficiency}
%%═══════════════════════════════════════════════════════════════════

By von Neumann's theorem
(\cite{ReedSimon1975}, Thm.~X.2),
$\Hspec$ is essentially self-adjoint iff
$n_\pm:=\dim\ker(\Hspec^*\mp i)=0$.

\begin{lemma}[Deficiency vectors are orthogonal
to $\{\Phi_n^{(\infty)}\}$]
\label{lem:deficiency-analytic}
Let $f\in\ker(\Hspec^*+i)$.
Then $\ip{f}{\Phi_n^{(\infty)}}=0$
for all $n\ge 0$.
\end{lemma}

\begin{proof}
\textit{Step~1.}
$\Phi_n^{(\infty)}\in\Dom(\Hspec)$
and
$\Hspec\Phi_n^{(\infty)}
=\lambda_n^{(\infty)}\Phi_n^{(\infty)}$
(Remark~\ref{rem:phi-in-domain}).

\textit{Step~2.}
Adjoint relation:
$\ip{\Hspec^*f}{\Phi_n^{(\infty)}}
=\ip{f}{\Hspec\Phi_n^{(\infty)}}$.

\textit{Step~3.}
Substitute
$\Hspec^*f=-if$
and
$\Hspec\Phi_n^{(\infty)}
=\lambda_n^{(\infty)}\Phi_n^{(\infty)}$:
$(\lambda_n^{(\infty)}+i)\ip{f}{\Phi_n^{(\infty)}}=0$.
Since $\lambda_n^{(\infty)}\in\mathbb{R}$,
$\lambda_n^{(\infty)}+i\ne 0$;
hence $\ip{f}{\Phi_n^{(\infty)}}=0$.
\end{proof}

\begin{remark}[Completeness $\Rightarrow$ ESSA]
\label{rem:deficiency-orthogonal}
$\ker(\Hspec^*\pm i)
\subseteq(\mathrm{span}\{\Phi_n^{(\infty)}\})^\perp$;
hence completeness $\Rightarrow$ $n_\pm=0$.
Converse not established without
a cyclic vector or spectral measure.
\end{remark}

\begin{hypothesis}[Completeness]
\label{hyp:decay}
$\{\Phi_n^{(\infty)}\}$ is complete
in $L^2(\mathbb{R})$.
\end{hypothesis}

%%═══════════════════════════════════════════════════════════════════
\section{Essential Self-Adjointness}
\label{sec:esa}
%%═══════════════════════════════════════════════════════════════════

\begin{theorem}[Deficiency indices]
\label{thm:deficiency}
Assume Hypothesis~\ref{hyp:decay}.
Then $n_\pm(\Hspec)=0$.
\end{theorem}

\begin{proof}
Any $f\in\ker(\Hspec^*+i)$ satisfies
$f\perp\Phi_n^{(\infty)}$ for all $n$
(Lemma~\ref{lem:deficiency-analytic});
Hypothesis~\ref{hyp:decay} gives $f=0$.
\end{proof}

\begin{corollary}[Essential self-adjointness]
\label{cor:esa}
Assume Hypothesis~\ref{hyp:decay}.
Then $\overline{\Hspec}=\Hspec^*$.
\end{corollary}

\begin{theorem}[Spectral identity and RH reduction]
\label{thm:sigma-identity}
Assume Hypothesis~\ref{hyp:decay}
and Hypothesis~CCM \cite{CCM2025}.
Then:
\begin{enumerate}[\rm(i)]
\item $\sigma_{\rm app}(\Hspec)=\sigma(\Hspec)$.
\item Every non-trivial zero $\rho$ of $\zeta$
      lies in $\sigma(\Hspec)$.
\item All non-trivial zeros satisfy
      $\mathrm{Re}(\rho)=\tfrac12$.
\end{enumerate}
\end{theorem}

\begin{proof}
(i)~ESSA gives $\sigma_{\rm app}=\sigma$
(\cite{ReedSimon1975}, Thm.~VII.13).
(ii)~Paper~VIII, Thm.~5.2 $+$ (i).
(iii)~CCM maps zeros to spectral parameters;
self-adjointness forces $\sigma(\Hspec)\subset\mathbb{R}$.
\end{proof}

\begin{remark}[Meta-corollary status]
\label{rem:meta-corollary}
Theorem~\ref{thm:sigma-identity}(iii)
is a \emph{conditional meta-corollary}:
it depends on
Hypothesis~\ref{hyp:decay} (open),
Hypothesis~CCM (conjectural),
and an injective spectral calibration
$\rho\mapsto\sigma(\Hspec)$
assumed via CCM but not independently verified.
This is not a proof of the Riemann Hypothesis.
\end{remark}

\begin{remark}[CCM status]
\label{rem:ccm-status}
Hypothesis~CCM \cite{CCM2025} is conjectural.
The converse (spectrum $\to$ zeros)
is not established.
\end{remark}

%%═══════════════════════════════════════════════════════════════════
\section{Open Problems}
\label{sec:problems}
%%═══════════════════════════════════════════════════════════════════

\begin{problem}[Identification]
\label{prob:identification}
Prove $\overline{\Hspec}=\HSOT$.
Both agree on $\mathrm{span}\{\Phi_n^{(\infty)}\}$;
under Hypothesis~\ref{hyp:decay}
and norm-resolvent convergence (Paper~X),
identification follows by density.
\end{problem}

\begin{problem}[Completeness]
\label{prob:completeness}
Prove Hypothesis~\ref{hyp:decay}.
Completeness $\Rightarrow$ $n_\pm=0$;
converse not established.
\end{problem}

\begin{problem}[Uniform alignment]
\label{prob:uniform-alignment}
Prove Hypothesis~\ref{hyp:uniform-alignment}.
Would follow from uniform gap bound
or resolvent control (Paper~X).
\end{problem}

\begin{problem}[Sharp rate]
\label{prob:opnorm-sharp}
Is $\|\Hc-\Hspec\|\asymp c^{-1/4}$?
\end{problem}

\begin{problem}[Unconditional ESSA]
\label{prob:esa-unconditional}
Prove ESSA without Hypothesis~\ref{hyp:decay}.
\end{problem}

%%─── Bibliography
\begin{thebibliography}{99}

\bibitem{PaperV}[Author],
\textit{Effective bandwidth and $L^2$-concentration},
Paper~V, 2025.

\bibitem{PaperVI}[Author],
\textit{Pointwise Mellin decay and $c^{-2/3}$ upper bound},
Paper~VI, 2025.

\bibitem{PaperVII}[Author],
\textit{Composite asymptotics and $c^{-1/2}$ upper bound},
Paper~VII, 2025.

\bibitem{PaperVIII}[Author],
\textit{Non-cancellation control, $c^{-1/2}$ lower bound,
and spectral connection to zeros of $\zeta(s)$},
Paper~VIII, 2026.

\bibitem{PaperX}[Author],
\textit{Mosco convergence and resolvent stability},
Paper~X, 2026.

\bibitem{PaperXI}[Author],
\textit{Spectral structure and density for the
limiting PSWF operator},
Paper~XI, 2026.

\bibitem{Slepian1961}
D.~Slepian and H.~O.~Pollak,
\textit{Prolate spheroidal wave functions,
Fourier analysis and uncertainty~I},
Bell System Tech.\ J.\ \textbf{40} (1961), 43--63.

\bibitem{Osipov2013}
A.~Osipov, V.~Rokhlin, H.~Xiao,
\textit{Prolate Spheroidal Wave Functions
of Order Zero},
Springer, 2013.

\bibitem{Olver1974}
F.~W.~J.~Olver,
\textit{Asymptotics and Special Functions},
Academic Press, 1974.

\bibitem{Kato1966}
T.~Kato,
\textit{Perturbation Theory for Linear Operators},
Springer, 1966.

\bibitem{ReedSimon1975}
M.~Reed and B.~Simon,
\textit{Methods of Modern Mathematical Physics,
Vol.~II},
Academic Press, 1975.

\bibitem{CCM2025}
A.~Connes, C.~Consani, H.~Moscovici,
\textit{Zeta Spectral Triples},
arXiv:2511.22755, 2025.

\end{thebibliography}

\end{document}
